\documentclass[11pt,a4paper]{article}
\usepackage[T1]{fontenc}
\usepackage[utf8]{inputenc}
\usepackage{lmodern,microtype,amsmath,amssymb,amsthm,booktabs,longtable,graphicx}
\usepackage[margin=23mm]{geometry}
\usepackage[hidelinks]{hyperref}
\usepackage{fancyhdr}
\pagestyle{fancy}\fancyhf{}\fancyhead[L]{FIN ST8548}
\fancyhead[R]{6 September 2026}\fancyfoot[C]{\thepage}
\setlength{\headheight}{14pt}\setlength{\emergencystretch}{2em}
\newtheorem{theorem}{Theorem}
\newtheorem{proposition}[theorem]{Proposition}
\newcommand{\one}{\mathbf1}
\newcommand{\arsinh}{\operatorname{arsinh}}
\newcommand{\diag}{\operatorname{diag}}
\hypersetup{pdftitle={FIN: Same Heat Flow, Different Fluctuation Actions},
 pdfauthor={Krzysztof Żuchowski}}
\begin{document}\hypersetup{pageanchor=false}
\begin{titlepage}\centering\vspace*{16mm}
{\Huge\bfseries Same FIN Heat Flow,\par}
\vspace{3mm}{\Huge\bfseries Different Fluctuation Actions\par}
\vspace{8mm}{\LARGE Cosh dissipation, integer-count constraints,\par}
{\LARGE and exact coarse current contraction\par}
\vspace{12mm}{\Large Programme ST8548\par}
\vspace{8mm}{\Large Krzysztof \.Zuchowski\par}
\vspace{3mm}{Independent Researcher --- Fractal Information Theory Project\par}
{ORCID: 0009-0002-0909-3613\par}
\vspace{9mm}{6 September 2026\par}
\vfill
\begin{minipage}{.87\textwidth}
An analytical and computational follow-up to ST8547. The report tests whether
the quadratic logarithmic-mean transport cost is a microscopic fluctuation
law. It derives an alternative from independent Markov walkers, identifies
observable differences, and proves exact current contraction through product
refinement. The microscopic realization remains a supplied model.
\end{minipage}
\vfill{\small English research preprint; CC BY 4.0.}
\end{titlepage}\hypersetup{pageanchor=true}

\begin{abstract}
The logarithmic mean gives a valid quadratic Onsager representation of
finite graph heat flow. Independent continuous-time Markov walkers with the
same conductances instead yield a nonquadratic fluctuation action with a
cosh dissipation potential. Both have the same mean current and the same
time-reversal cost asymmetry. Their local variances differ by
\(z\coth z\ge1\), where \(z\) is half the log population ratio.
For the explicit strict FIN test the ratio is \(1.4683539294\).
The quadratic prescription even violates the variance lower bound for
integer-valued Markov counts when \(|\log(a/b)|>2\).

The exact jump action contracts under summing fine channels for arbitrary
positive hidden conditional distributions. Its covariance also projects
exactly. This differs from the quadratic Onsager equality criterion in
ST8547, which requires independence. Thus a defect of one gradient metric
does not imply failure of exact stochastic coarse-graining.

All conclusions are conditional on supplied rates and a microscopic counting
model. The underlying nonquadratic gradient structure is established
literature. This study provides a checked application, a counting test and
corrections to FIN's interpretation; it makes no claim of new fundamental
physics or mathematical priority.
\end{abstract}

\tableofcontents\newpage

\section{Scope, source evidence and conventions}

The project state map in AGENTS.md, the ST8547 proofs and code, and the
earlier configuration-generator examples were consulted. This is a focused
proof-level continuation, not a new verification of every historical FIN
report. The raw archived status labels and hashes are not treated as
independent proof certificates.

The basic microscopic/gradient distinction was established by Mielke,
Peletier and Renger \cite{mpr}, with further developments by Liero et al.
\cite{lmpr}. Maas \cite{maas} supplies the logarithmic-mean quadratic
representation. The present derivations and numerical tests state exactly
which of these structures is being used.

The strict conductances are
\[
W_{ij}=\frac{\cos(0.18575d+0.16250)}{1+d^{9/5}},\qquad
d=\min(|i-j|,12-|i-j|),\quad W_{ii}=0.
\]
They are positive on \(d=1,\ldots,6\). Let
\(Q=W-\diag(W\one)=-A\), a row-generator convention.
The raw legacy kernel has signed weights and is not substituted into this
positive-rate construction.

Assume \(N\) independent Markov walkers with generator \(Q\). Their empirical
occupation vector has increments \((e_j-e_i)/N\), occurring at rate
\(Np_iW_{ij}\). Its exact generator is
\begin{equation}
(\mathcal G_NF)(p)=N\sum_{i\ne j}p_iW_{ij}
\left[F\!\left(p+\frac{e_j-e_i}{N}\right)-F(p)\right]. \label{empirical}
\end{equation}
Coordinate functions give the mean equation \(\dot p=-Ap\).
The population, independence assumption and rate interpretation of \(W\)
are additional model premises; physical time is not derived.

\section{Exact local current fluctuations}

Orient edge \(i\to j\). Write
\[
a=W_{ij}p_i,\qquad b=W_{ij}p_j,\qquad a,b>0.
\]
The local current Hamiltonian is
\begin{equation}
h(s;a,b)=a(e^s-1)+b(e^{-s}-1). \label{hamiltonian}
\end{equation}
It follows either from the two oriented generator increments or from the
nonlinear exponential transform of \eqref{empirical}. Derivatives at zero
give instantaneous cumulant rates
\[
\kappa_1=a-b,\quad\kappa_2=a+b,\quad
\kappa_3=a-b,\quad\kappa_4=a+b.
\]
For \(N\) independent walkers the total-count cumulants scale by \(N\).

This is an exact \emph{local} action density. For a closed chain it is not
the finite-time count distribution at arbitrary time: the state changes
and return transitions occur. Section~\ref{verification} checks finite-time
counts with the full tilted generator.

\begin{theorem}[Current rate function]
The Legendre transform of \eqref{hamiltonian} is
\begin{equation}
\ell(j;a,b)=j\arsinh\frac{j}{2\sqrt{ab}}
-\frac j2\log(a/b)-\sqrt{j^2+4ab}+a+b. \label{cost}
\end{equation}
It is strictly convex in \(j\), with unique zero \(j_0=a-b\).
\end{theorem}
\begin{proof}
The maximizing tilt solves \(j=ae^s-be^{-s}\), so
\[
s_*=\arsinh\frac{j}{2\sqrt{ab}}-\frac12\log(a/b).
\]
At that tilt \(ae^s+be^{-s}=\sqrt{j^2+4ab}\), yielding \eqref{cost}.
Its second derivative is \(1/\sqrt{j^2+4ab}>0\), and direct substitution
at \(a-b\) gives zero.
\end{proof}

There is an independent flux representation:
\[
\ell(j;a,b)=
\inf_{u,v\ge0:\,u-v=j}
\left[u\log(u/a)-u+a+v\log(v/b)-v+b\right].
\]
The minimizers are
\(u=(\sqrt{j^2+4ab}+j)/2\),
\(v=(\sqrt{j^2+4ab}-j)/2\).
Relative-entropy perspective convexity proves joint convexity in
\((j,a,b)\). Also
\(\ell(cj;ca,cb)=c\ell(j;a,b)\).

\section{The nonquadratic entropy-dissipation structure}

Set \(g=2\sqrt{ab}\), \(z=\frac12\log(a/b)\), and define
\[
\mathcal R(j)=j\arsinh(j/g)-\sqrt{j^2+g^2}+g,\qquad
\mathcal R^*(\zeta)=g(\cosh\zeta-1).
\]
These are symmetric Legendre-dual potentials, and
\[
\ell(j;a,b)=\mathcal R(j)+\mathcal R^*(z)-jz,\qquad
j_0=\partial_\zeta\mathcal R^*(z)=a-b.
\]
The force \(z\) corresponds to one-half of negative relative Shannon
entropy in the energy-dissipation convention. This factor follows from
the reversal identity; it is not a physical constant.

The same microscopic independence premise also gives relative entropy from
equilibrium counting. For invariant law \(\pi\), independent occupations
are multinomial:
\[
\Pr(Np)=\frac{N!}{\prod_i(Np_i)!}\prod_i\pi_i^{Np_i}.
\]
Stirling's formula yields
\(-N^{-1}\log\Pr(Np)\to D(p\Vert\pi)\).
Here \(\pi_i=1/12\) since \(W\) is symmetric.
Thus a microscopic counting assumption can select a conditional entropy
and dissipation structure together. It is not furnished by \(A\) alone.

\section{Quadratic mobility does not fix microscopic noise}

Let \(m=\Lambda(a,b)=(a-b)/\log(a/b)\), with diagonal value \(m=a\).
To compare fairly, normalize the quadratic action as
\begin{equation}
\ell_{\rm Q}(j;a,b)=\frac{[j-(a-b)]^2}{4m}. \label{quadratic}
\end{equation}
It has the same mean and exactly the same reversal asymmetry:
\[
\ell(j)-\ell(-j)=\ell_{\rm Q}(j)-\ell_{\rm Q}(-j)
=-j\log(a/b).
\]
The factor 4 compensates the one-half entropy convention; it is chosen
explicitly rather than hidden in a clock.

\begin{theorem}[Variance mismatch]
The jump action and the quadratic prescription have local variances
\[
v_{\rm jump}=a+b,\qquad v_{\rm Q}=2\Lambda(a,b).
\]
Their ratio is
\[
\boxed{\frac{v_{\rm jump}}{v_{\rm Q}}=z\coth z\ge1,}
\]
with equality precisely at \(a=b\).
\end{theorem}
\begin{proof}
Inverse rate-function curvature at \(j_0\) equals \(a+b\) for
\eqref{cost} and \(2m\) for \eqref{quadratic}. Write
\(a=\sqrt{ab}e^z,b=\sqrt{ab}e^{-z}\).
The ratio becomes \(z/\tanh z\), and
\(|\tanh z|\le|z|\), strictly off zero.
\end{proof}

Near equilibrium this ratio is \(1+z^2/3+O(z^4)\).
Even at equilibrium, the exact fourth cumulant rate is \(2a>0\), while
the local quadratic Hamiltonian has zero fourth cumulant.
The deterministic ST8547 gradient identity remains valid.

\begin{theorem}[Integer-count consistency]
For any Markov additive count with integer increments \(k\), nonnegative
local rates \(c_k\), and finite first two moments,
\[
v=\sum_k k^2c_k\ge\left|\sum_k kc_k\right|=|\mu|.
\]
The quadratic prescription violates this necessary condition whenever
\[
\boxed{|\log(a/b)|>2.}
\]
\end{theorem}
\begin{proof}
Use \(k^2\ge|k|\) and the triangle inequality.
For \eqref{quadratic},
\(v_{\rm Q}/|j_0|=2/|\log(a/b)|\), which is less than one in the stated
regime.
\end{proof}

This excludes an exact integer-valued Markov count interpretation in that
regime, even allowing collective integer jumps. It does not exclude a
continuous Gaussian approximation or another coarse observable.

Three means have distinct roles: the logarithmic mean relates mean current
to finite entropy force; the geometric mean sets the cosh prefactor; the
arithmetic combination determines local shot noise.

\subsection{Whole-graph covariance}

The empirical-occupation covariance rate per particle is
\[
\Sigma(p)=\sum_{i<j}W_{ij}(p_i+p_j)
(e_i-e_j)(e_i-e_j)^\top.
\]
The quadratic prescription uses \(2K_\Lambda(p)\).
Arithmetic-logarithmic mean ordering yields
\[
\boxed{\Sigma(p)-2K_\Lambda(p)\succeq0.}
\]
The difference is a graph Laplacian with nonnegative conductances.
For connected support it is zero exactly at uniform \(p\).

\section{Exact coarse current contraction}

Let one coarse edge aggregate fine channels with positive rates
\(a_b,b_b\). Set \(a=\sum_ba_b,b=\sum_bb_b\).

\begin{theorem}[Current-action contraction]
\[
\boxed{\inf_{\sum_bj_b=j}\sum_b\ell(j_b;a_b,b_b)=\ell(j;a,b).}
\]
No equality of forward and reverse conditional child profiles is required.
\end{theorem}
\begin{proof}
Choose \(s\) with \(ae^s-be^{-s}=j\), and
\(j_b=a_be^s-b_be^{-s}\). All fine costs have derivative \(s\) at this
feasible vector, so convexity proves optimality. Their Legendre values
sum to \(sj-\sum_bh(s;a_b,b_b)=sj-h(s;a,b)\).
\end{proof}

In the supplied product FIN graph, \(a_b=W_{ij}P_{ib}\) and
\(b_b=W_{ij}P_{jb}\). Consequently this contraction holds for arbitrary
strictly positive correlated fine states. Its covariance also satisfies
\[
C\Sigma_{\rm fine}(P)C^\top=\Sigma_{\rm coarse}(CP),
\]
since the edge coefficients are linear in \(P\) and vertical increments
vanish after projection.

For \(A\otimes I+I\otimes B\), the base-coordinate Markov process is exactly
lumpable for every initial joint state. This is a product generator, not
the interacting Ising/prism generator considered in older stationary-RG
studies. Their different lumpability properties are not contradictory.

The quadratic contraction instead gives
\[
\inf_{\sum j_b=j}\sum_b\ell_{\rm Q}(j_b;a_b,b_b)
=\frac{[j-(a-b)]^2}{4\sum_b\Lambda(a_b,b_b)}.
\]
As shown in ST8547,
\(\sum_b\Lambda(a_b,b_b)\le\Lambda(a,b)\).
For \(j\ne j_0\), equality of quadratic costs requires matching forward
and reverse normalized child profiles.

Thus the ST8547 correlation defect is a defect of that quadratic
representation; it does not force a defect of exact Markov counting
statistics. Its KL entropy-production identity remains valid.

\section{Strict numerical example and independent checks}\label{verification}

Use \(p_i=(i+1)/78\) and edge \(0\to11\).
\begin{center}
\begin{tabular}{lr}
\toprule
Quantity & Value\\
\midrule
\(W_{0,11}\) & \(0.4699856726450201\)\\
\(a\) & \(0.006025457341602822\)\\
\(b\) & \(0.07230548809923387\)\\
Mean current \(a-b\) & \(-0.06628003075763104\)\\
Jump variance rate & \(0.0783309454408367\)\\
Quadratic variance rate & \(0.05334609311241991\)\\
Jump/quadratic variance & \(1.4683539294201824\)\\
\bottomrule
\end{tabular}
\end{center}
The population ratio is 12. The local count ratios are
\[
\frac{v_{\rm jump}}{|j_0|}=\frac{13}{11}>1,\qquad
\frac{v_{\rm Q}}{|j_0|}=\frac2{\log12}<1.
\]

\begin{figure}[ht]
\centering\includegraphics[width=\textwidth]{findings.pdf}
\caption{The exact and quadratic current costs share the deterministic
minimum (left). Their local variance-to-mean ratios differ, and the
quadratic prescription crosses below the integer-jump bound (right).}
\end{figure}

For forward child proportions \((0.8,0.2)\), reverse proportions
\((0.25,0.75)\) and current \(j=0.04\), the exact contraction error is
\(2.78\times10^{-17}\). The quadratic fine-contracted and coarse costs are
\(0.1178653956\) and \(0.1058694675\).
The graph covariance projection error is \(1.11\times10^{-16}\);
the quadratic projection defect norm is \(0.0079424221\).

Finite-time counts were independently computed using derivatives of the full
tilted 12-state generator. For one walker with initial distribution \(p\):
\begin{center}
\begin{tabular}{rrr}
\toprule
Window \(t\) & variance/\(t\) & fourth cumulant/\(t\)\\
\midrule
0.1 & 0.0745294243 & 0.0717795571\\
0.01 & 0.0779222577 & 0.0776108767\\
0.0001 & 0.0783268250 & 0.0783236669\\
0.00001 & 0.0783305334 & 0.0783302175\\
\bottomrule
\end{tabular}
\end{center}
The limits are the local jump value \(0.0783309454\).
No finite-time Poisson approximation was substituted for this calculation.

Eighteen tests independently check the Legendre transform by optimization,
the flux-relative-entropy contraction, high-precision derivatives, the
cosh decomposition, convexity, integer-count bound, covariance order,
projection, finite-time tilted moments and the correction below.
These are numerical checks supporting explicit analytical proofs,
not interval certificates or physical evidence.

\section{A necessary correction to the older kinetic example}

In ST2547--ST2561, the three-bit model has energy
\(E=-\theta x_1x_2x_3\). Every flip has
\(|\Delta E|=2|\theta|\). Therefore
\[
Q_{\rm Met}=(1+e^{-2|\theta|})Q_{\rm HB},\qquad
Q_{\rm sym}=2\cosh(\theta)Q_{\rm HB}.
\]
They have identical normalized spectral shapes. The earlier claim that
this example distinguished more than a clock factor is false.
At \(\theta=0.3\), both proportionality identities replay within
\(4.45\times10^{-16}\).

General kinetic nonuniqueness remains, but that specific example did not
prove it beyond time scaling. Historical files are preserved and this
correction supersedes their interpretation.

\section{Remaining assumptions and the next physical test}

Microscopic independence is essential. If particles instead move in
synchronized groups of size \(s\), the local Hamiltonian per total particle is
\[
h_s(\xi)=\frac1s\left[a(e^{s\xi}-1)+b(e^{-s\xi}-1)\right].
\]
Its mean remains \(a-b\), but its variance is \(s(a+b)\). Hence the same
mean heat equation does not determine a unique population law. The
12-state walker model is also distinct from a 4096-state binary
configuration model.

A dimensionless conditional experiment can compare short-window net-count
variance to the absolute mean net count. Population size and the common clock
factor cancel. At zero drift, compare local fourth and second cumulants
instead. Repeated preparation, bidirectional event counts, finite-window
extrapolation and detector-loss checks are needed; no such laboratory record
is claimed here.

The main result for FIN is an observable distinction between gradient
representations with identical deterministic dynamics, and exact coarse
contraction of the microscopic current action. The best next investigation
is a finite-noise counting protocol testing single versus collective events
and the predicted cumulants. A sourced FIN microscopic law must specify what
those events are.

The strict \(U_t=e^{-itA}\) and \(P_t=e^{-tA}\) remain distinct functional
calculi of one operator. This study concerns the stochastic realization of
the heat channel. It derives neither quantum measurement dynamics nor
physical units, selector/QW-2191, legacy role transfer, continuum spacetime,
Standard Model, gravity, \(L_{\rm total}\), or a Theory of Everything.

\appendix
\section{Reproduction}
\begin{verbatim}
python3 analysis.py --output replay_results.json
python3 -m unittest -v test_analysis.py
MPLCONFIGDIR=/tmp/fin-st8548-mpl python3 plot_results.py
pdflatex -interaction=nonstopmode -halt-on-error report.tex
pdflatex -interaction=nonstopmode -halt-on-error report.tex
\end{verbatim}
The package includes the full proofs in PROOFS\_EN.md, input hashes and
dependency versions in provenance.json, and checksums. No network access or
large computation is required for replay.

\begin{thebibliography}{9}
\bibitem{maas} J. Maas, \emph{Gradient flows of the entropy for finite Markov
chains} (2011). \url{https://arxiv.org/abs/1102.5238}.
\bibitem{mpr} A. Mielke, M. A. Peletier and D. R. M. Renger,
\emph{On the relation between gradient flows and the large-deviation principle,
with applications to Markov chains and diffusion}.
\url{https://arxiv.org/abs/1312.7591}. See Section 4.1 and Example 4.2.
\bibitem{lmpr} M. Liero, A. Mielke, M. A. Peletier and D. R. M. Renger,
\emph{On microscopic origins of generalized gradient structures},
Discrete and Continuous Dynamical Systems--S 10 (2017), 1--35.
\url{https://doi.org/10.3934/dcdss.2017001}.
\end{thebibliography}
\end{document}
